% Page 087

\seite{87}

\margmark{1. Juli}%
In den Tagen vom 14.\ bis 18.\ Juni wurden\ob
die beiden unteren Fenster, welche außer noch\ob
einige Reparaturen die Geistliche vorgenommen\ob
an werden waren, aufgehängte Glockenring\ob
wird ein brandfreies verlängertes großes\ob
Fenster, jetzt am Eingang hinten die Kapelle ge=\ob
schüttet werden können, worauf die Glockentuben\ob
für ein anderes Conto {[gesezt]}. Diese Anlage ist\ob
fast zu begrüßen. Da der Vorsteher des Ortes\ob
ch {[grossgewöhnter]} Weise das Lasten versucht jetzt\ob
gemeinsam wird. Und diese Verbesserung ist\ob
die einmüthigsten Vorrath des Kapellen= Küster\ob
Peter Manns zu verdanken, der das Projekt ge=\ob
gänzlich vorbereitet und nicht eher, bis es\ob
verwirklicht war. Manns hat sich somit einen\ob
Namen in der Geschichte der Kapelle erworben.

Mit dem heutigen Tag beginnen die\ob
großen Dir Ferien bis einschließlich\ob
14.\ Juli.

\margmark{15. August}%
Der Sommer gestaltet sich zur vollen\ob
Zufriedenheit der Landwirthschaft. Die Heu=\ob
ernte ist fast nicht ausgefallen. Das Ernte des\ob
Getreidesaate ist ein einziges und dazu\ob
fast vollständig. Da allmählich eine bessere\ob
Friede begann und wenig Grünfutter für\ob
viel Regen vorgefallen sind, droßte es nur\ob
wenig Zwiebel Futter geben.

\margmark{20. Septembr}%
Die Herbstferien dauern vom 18.\ Sep=\ob
tember bis 22.\ Oktober einschließlich

Nach den ersten Schnipsen zu ur=\ob
theilen gibt es lang eine schlechte Kartoffel=\ob
ernte, doch die großen Vorarbeit der Kar=\ob
toffelernte. Die ersten Preise pro 36.\ Pfund 25--55.

\margmark{1. Novembr}%
Gleichzeitig ist der Landmann mit\ob
seiner Arbeit, bezüglich eines guten Weiten\ob
fertig geworden. Er hat überall mal ein gutes\ob
Jahr hinter sich. Er ist reichlich versorgt und zufrieden.
\pend
